\chapter{\glsfmtfull{xss}}
\glsxtrlong{xss} attacks are a type of injection in which malicious scripts are injected into trusted websites. These attacks occur when a web application takes untrusted data and sends it to a web browser without proper validation or escaping, generally in the form of a browser-side script that rewrites the content of an HTML page. 

\begin{description}
    \item[Severity] Very Severe
    \item[Priority] High
    \item[Prevalence] Frequently discovered
    \item[Scope] All web applications and user browsers
    \item[Technical Impact] Execution of unauthorized scripts in the user's session context
    \item[Worst-Case Consequences] Theft of session tokens, sensitive data extraction, and client-side system compromise
\end{description}

\noindent \glsxtrlong{xss} is incredibly dangerous because the end user's browser has no way to know that the executed script should not be trusted. Because the browser thinks the script came from a trusted source, the malicious script can access any cookies, session tokens, or other sensitive information retained by the browser and used with that site.

\paragraph{Underlying Idea}
Let us assume our application uses the following PHP script to capture a search query and display it back to the user:

\begin{lstlisting}[language=PHP]
<?php
// http://search.com/?query=Q-bits
$query = $_GET["query"];

if (isset($query)) {
    echo "Search results for: " . $query;
}
?>
\end{lstlisting}

Normally, an \gls{http} request is performed by the client querying something benign like \texttt{Q-bits}. The server reads the parameter and returns \texttt{Search results for: Q-bits}.

However, a malicious user can craft a URL containing active HTML or JavaScript instead of normal text. For example, if the attacker inputs \texttt{<script>alert(document.cookie);</script>}, the response generated by the server becomes:

\begin{lstlisting}[language=HTML]
<html>
Search results for: 
<script>alert(document.cookie);</script>
</html>
\end{lstlisting}

The result is a direct attack on the user interacting with the page. The browser executes the injected script, potentially sending the user's active session cookie to an attacker-controlled server.

\paragraph{Flow and Steps} 
XSS attacks typically follow a standard flow:
\begin{enumerate}
    \item The attacker identifies a website with \gls{xss} vulnerabilities (e.g., a banking portal).
    \item The attacker creates a malicious URL or script that utilizes the vulnerability.
    \item The attacker uses social engineering to trick the victim into clicking the malicious URL or executing the script.
    \item The victim clicks the URL, sending a request enriched with malicious code to the vulnerable server.
    \item The vulnerable website builds a response page for the victim that includes the unescaped malicious code.
    \item The victim's browser receives the response and executes the malicious payload.
\end{enumerate}

\paragraph{Examples in Different Languages}
The core mechanism behind an \glsxtrshort{xss} vulnerability is that untrusted user input is directly output to the client web page without any sanitization or encoding.

\textbf{PHP}
\begin{lstlisting}[language=PHP]
<?php
// Unintended use: ?group=window.location="https://maliciouswebsite.com"
echo "<p>Hello user, your current group is [" . $_GET['group'] . "]</p>";
?>
\end{lstlisting}

\textbf{Java (Spring Servlet)}
\begin{lstlisting}[language=Java]
@GetMapping("/direct")
public void directLink (@RequestParam String param, HttpServletResponse response) throws IOException {
    // ...
    var writer = response.getWriter();
    // Vulnerable string concatenation without HTML encoding
    writer.write("<div class=\"panel-heading\"><h1>"+ param + "</h1></div>");
    response.getWriter().flush();
}
\end{lstlisting}

These vulnerabilities are not restricted to basic HTML tags. Attackers can leverage multiple attack vectors including the \texttt{<body>} tag (via \texttt{onload} attributes), \texttt{<img>} tags (via malicious \texttt{src} or \texttt{onerror} handlers), \texttt{<iframe>} tags, and even \texttt{<link>} or CSS \texttt{style} attributes.

\section{Different Types of \glsfmtshort{xss}} 
\gls{xss} vulnerabilities are classified by how the malicious payload is delivered and executed within the application environment.

\begin{itemize}
    \item \textbf{Type II: Reflected (Non-Persistent) \gls{xss}}: User input is immediately returned by a web application in an error message, search result, or any other response without being made safe to render. 
    \begin{itemize}
        \item The user-provided data is not permanently stored on the server.
        \item It relies heavily on social engineering, such as SPAM emails with crafted links, or malicious posts on social media.
    \end{itemize}
    
    \item \textbf{Type I: Stored (Persistent) \gls{xss}}: User input is collected from the vulnerable application and permanently stored on the target server (e.g., in a database, message forum, or visitor log).
    \begin{itemize}
        \item A victim is attacked simply by browsing the application and retrieving the stored data, without requiring a direct social engineering link.
        \item Attackers inject malicious code into the stored data (like a forum post), and the payload executes every time a user views that specific post.
    \end{itemize}
    
    \item \textbf{Type 0: DOM-Based \gls{xss}}: A form of \gls{xss} where the entire tainted data flow from source to sink takes place entirely in the browser.
    \begin{itemize}
        \item The source of the data is in the Document Object Model (DOM) (e.g., \texttt{document.location.href}), and the sink is also in the DOM (e.g., \texttt{document.write} or \texttt{innerHTML}).
        \item The malicious data flow never leaves the browser, meaning server-side attack detection tools may fail to see or detect the attack.
    \end{itemize}
\end{itemize}

\section{\glsfmtshort{xss} Mitigation}
No single technique can completely solve \gls{xss}; a combination of strategies must be implemented.

\begin{enumerate}
    \item \textbf{Don't Trust User Input (Effectiveness: \textcolor{orange}{MEDIUM})} \\
    Treat all user input as untrusted, including data from authenticated or internal users.
    
    \item \textbf{Set the HttpOnly Flag (Effectiveness: \textcolor{orange}{MEDIUM})} \\
    Configure the HTTP response header to include the \texttt{HttpOnly} flag on sensitive cookies. This ensures such cookies cannot be accessed or stolen via client-side JavaScript.
    
    \item \textbf{Content Security Policy (CSP) (Effectiveness: \textcolor{olive}{HIGH-})} \\
    Implement a CSP via HTTP response headers to declare a strict whitelist of allowed dynamic resources. This defines exactly which scripts or remote links are permitted to execute in the browser.
    
    \item \textbf{HTML Sanitization (Effectiveness: \textcolor{green!60!black}{HIGH})} \\
    If user input absolutely needs to contain HTML markup (e.g., a rich text editor), use a trusted and verified library to parse and clean the HTML, stripping out any dangerous tags or event handlers.
    
    \item \textbf{Output Escaping/Encoding (Effectiveness: \textcolor{green!60!black}{HIGH})} \\
    This is the primary defense. Ensure that any user input is properly encoded (e.g., HTML, JavaScript, CSS, or URL encoding) before it reaches the output statement, forcing the browser to render it safely as text rather than executable markup.
\end{enumerate}

\paragraph{Mitigation via Encoding}
Modern web development relies on dedicated encoding libraries to neutralize characters that have special meaning in HTML (such as \texttt{<}, \texttt{>}, \texttt{\&}, \texttt{"}, and \texttt{'}).

\textbf{PHP (Using htmlspecialchars)}
\begin{lstlisting}[language=PHP]
<?php
$query = $_GET['query'];
if (isset($query)) {
    // Converts special characters to HTML entities (e.g., < becomes &lt;)
    // This prevents the execution of tags on the user's browser
    $query = htmlspecialchars($query); 
    echo "Search results for: $query";
}
?>
\end{lstlisting}

\textbf{Java (Using Apache Commons Text or OWASP Encoder)}
\begin{lstlisting}[language=Java]
// Input validation/sanitization is required before writing it to the response.
// Use libraries like StringEscapeUtils.escapeHtml4() or Encode.forHtmlContent()
var writer = response.getWriter();

writer.write("<div class=\"panel-heading\"><h1>"+ 
             StringEscapeUtils.escapeHtml4(param) + "</h1></div>");

String output = "<div class=\"panel-body\"> <ul><li>%s</li></ul> </div>";

writer.write(String.format(output,
             StringEscapeUtils.escapeHtml4(prod.getDescription())));
\end{lstlisting}

\section{Tools and Affected Technologies}
\gls{xss} is an incredibly pervasive vulnerability because it exploits the trust between the browser and the application's output, rather than exploiting a specific programming language. Any programming language used to build a website—including PHP, ASP, C\#, VB.Net, Java, Perl, and CGI—is vulnerable if it fails to escape untrusted input before rendering it to the client.

\section{Chapter Summary}

\begin{concept}{Core \glsfmtshort{xss} Concepts}{xss-summary}
    \textbf{The Fundamental Problem:} User input is directly displayed in an output web page without appropriate sanitization or encoding, allowing attackers to inject client-side scripts.
    
    \textbf{Primary Defenses:}
    \begin{itemize}
        \item \textbf{Escape / Encode Output:} Encode the output using HTML encoding so that any malicious link or JS code remains uninterpreted by the browser.
        \item \textbf{Sanitize Input:} Sanitize any user input which might reach output statements, including statements that write to a database or save cookies.
        \item \textbf{Defense in Depth:} Combine encoding with Content Security Policies (CSP) and \texttt{HttpOnly} cookies to limit the damage if a vulnerability is ever introduced.
    \end{itemize}
\end{concept}